\section{Apple: resources available for payment}
\label{sec:apple-liquidity}

The procurement and settlement choices just described determine when Apple
needs money. Funding those payments requires a second account: what resources
can be used, how securities become cash, and when customers pay. Annual and
interim Note 4 distinguish these questions even though a single total might
appear to answer them. The following table combines the notes' measurement
and balance-sheet classifications. Amounts are in millions of US dollars at
the displayed precision; the December date is the end of the first fiscal
quarter, not another annual endpoint
\citep[Note 4, printed pp.~36--37]{apple2024}
\citep[Note 4, printed pp.~7--8]{apple2025q1}.

\begin{center}
\begin{tabular}{lrrr}
\toprule
Resource / measurement &30 Sep 2023&28 Sep 2024&28 Dec 2024\\\midrule
Cash and cash equivalents &29,965&29,943&30,299\\
Current marketable securities &31,590&35,228&23,476\\
Non-current marketable securities &100,544&91,479&87,593\\\midrule
Total fair value &162,099&156,650&141,368\\
Total adjusted cost &173,752&160,600&147,088\\
Net unrealised gains / losses &$-11,653$&$-3,950$&$-5,720$\\\bottomrule
\end{tabular}
\end{center}

Cash plus the two securities rows equals total fair value. Adjusted cost plus
net unrealised gains and losses provides a second reconciliation of that same
total; it is not an additional resource to add to it. The source calls the
column \emph{adjusted cost}. Relabelling every investment as amortised-cost
debt would erase the different instruments represented in the table.

\subsection{Restriction changes what can be spent}

At September 2024, cash includes 2.6 billion dollars held in escrow, while
marketable securities include a further 13.2 billion. Both are restricted from
general use and designated for the State Aid tax obligation discussed earlier.
Subtracting their displayed centres from the total gives
$156{,}650-2{,}600-13{,}200=140{,}850$ million. MD\&A separately reports
unrestricted cash, equivalents and securities of 140.8 billion
\citep[printed pp.~25--26 and Note 4]{apple2024}. The apparent difference
does not establish conflicting disclosures: the restrictions and the MD\&A
total are rounded to tenths of a billion, whereas the table is in millions.
If each amount is rounded to its nearest displayed quantum, the subtraction
has possible bounds of 140,749.5 and 140,950.5 million. Those bounds overlap
the 140,750 to 140,850 million rounding range for the MD\&A total. This is a
conditional compatibility check, not an issuer-stated rounding convention
or a confidence interval.

The December footnote makes the distinction between maturity and restriction
especially clear: \emph{current} marketable securities include 3.3 billion
dollars of restricted escrow funds, and cash includes 126 million. At the
displayed centres, total resources less this disclosed escrow are
$141{,}368-3{,}300-126=137{,}942$ million. Cash plus current securities is
53,775 million; removing the same escrow gives 50,349 million. These are
indicative subtractions, not exact freely spendable balances. Under the same
nearest-rounding hypothesis the first subtraction ranges from 137,891 to
137,993 million. Interim MD\&A supplies no numerical unrestricted total
\citep[Note 4 and printed p.~16]{apple2025q1}. A current classification
therefore cannot serve as permission to spend an asset, and excluding escrow
does not establish immediate liquidation at the reported price.

Maturity supplies a different kind of timing information. Note 1 tests the
three-month cash-equivalent maturity at \emph{purchase}, while annual Note 4
classifies marketable debt as available-for-sale and assigns current or
non-current status using each instrument's underlying maturity
\citep[printed pp.~34 and 37]{apple2024}. Among non-current debt excluding
mortgage- and asset-backed securities, the September maturity shares are
86\%, 10\% and 4\% for 1--5, 5--10 and more than 10 years. The separate
mortgage- and asset-backed population has shares of 14\%, 9\% and 77\%.
At December these become 85/11/4\% and 13/10/77\%, respectively
\citep[Note 4, printed p.~8]{apple2025q1}. Each triple divides its own
population; the two triples cannot be added. Nor do contractual maturity
bands establish a sale date, sale price or cash receipt schedule. Management's
stated belief in sufficient resources also relies on ongoing operating cash
and continued access to debt markets. It is a dated funding assessment,
rather than a solvency guarantee derived from the table.

\subsection{Security turnover differs from valuation}

Securities can fund payments through sales or maturities, but a fall in their
reported balance need not equal cash released. Fiscal 2024 cash proceeds from
maturities and sales are 51,211 and 11,135 million, while purchases use 48,656.
Net cash released is therefore $51{,}211+11{,}135-48{,}656=13{,}690$ million.
The marketable-security balance falls by only
$126{,}707-132{,}134=-5{,}427$ million
\citep[cash-flow statement, printed p.~33, and Note 4]{apple2024}.
To expose the missing reconciliation, let $M_0$ and $M_1$ be total current
and non-current marketable securities at the period's opening and closing.
Let $C$ be maturity and sale proceeds less purchases, positive when cash is
released. Define $R$ as the unallocated difference between the stock movement
and that cash measure, all in millions over the same period:
\[
\begin{aligned}
 M_1-M_0&=-C+R,\\
 \text{Fiscal 2024:}\quad -5{,}427&=-13{,}690+8{,}263,\\
 \text{First quarter 2025:}\quad -15{,}638&=-13{,}335-2{,}303.
\end{aligned}
\]
For the quarter, the cash measure is
$15{,}967+3{,}492-6{,}124=13{,}335$ million
\citep[cash-flow statement and Note 4]{apple2025q1}. The relation defines a
residual; it does not explain it. Cash sale proceeds differ from the carrying
amount removed on disposal, and changes in measurement or classification may
also enter the stock. Apple's specific-identification policy for securities
sold requires the disposed lots to recover their cost. Neither the positive
8,263 nor the negative 2,303 identifies a fair-value, currency or
reclassification journal.

The valuation columns help separate another tempting explanation. Across the
annual dates, total adjusted cost falls by 13,152 million and net unrealised
gains and losses rise by 7,703, producing a total fair-value decline of 5,449.
Across the quarter, the corresponding changes are $-13{,}512$ and
$-1{,}770$, producing $-15{,}282$. These totals include cash and equivalents,
whereas $M_0$ and $M_1$ include only marketable securities. They also differ from the
comprehensive-income statement's net-of-tax marketable-debt amounts: 6,054
million for the year and $-1{,}427$ for the quarter, including reclassification
adjustments \citep[comprehensive-income statement, printed p.~31]{apple2024}
\citep[comprehensive-income statement, printed p.~2]{apple2025q1}. Endpoint
unrealised balances and period OCI differ in population, tax and timing;
their numerical differences cannot be assigned to a tax charge or realised
gain without the missing movements.

Even reported cash needs matching populations. Annual cash-flow cash, cash
equivalents and restricted cash opens at 30,737 million, exceeding opening
balance-sheet cash and equivalents of 29,965 by 772. Both close at 29,943.
Thus the cash-flow change is $-794$, while the balance-sheet line changes by
$-22$; the cause and location of the earlier 772 remain unallocated. The
quarter's endpoint populations agree, giving the same 356 million increase
\citep[balance and cash-flow statements]{apple2024,apple2025q1}.

Valuation risk is likewise distinct from a payment. The notes describe Level 2
prices or models with significant observable inputs and counterparties that
\emph{generally} have high credit ratings. Annual Item 7A reports that a
hypothetical 100-basis-point rate rise at all tenors would reduce investment
fair value, including associated derivatives, by 2,755 million at September
2024 and 3,089 million at September 2023 \citep[printed p.~27]{apple2024}.
Those are dated issuer scenarios, not cash outflows or loss probabilities;
the issuer's valuation model has not been replicated. Interim Item 3's
no-material-change referral does not issue a new December sensitivity.

\subsection{Future receipts need customers and collection dates}

Receivables may turn into cash later, but their movement does not identify
collections. Annual net trade receivables rise by 3,902 million, while the
indirect cash-flow adjustment is $-3{,}788$. Their sum leaves an unallocated
114 million difference. In the quarter, the balance change of $-3{,}771$
and adjustment of $+3{,}597$ leave $-174$. Vendor non-trade changes match
the opposite indirect adjustments exactly: $+1{,}356/-1{,}356$ annually
and $-3{,}166/+3{,}166$ quarterly. Even that match supplies neither gross
sales nor total collections \citep[balance and cash-flow statements]{apple2024,apple2025q1}.
Sales, allowances, currency and noncash movements still need their own
evidence before receipts can be dated.

Concentration describes who might pay, but overlapping descriptions cannot be
summed as separate exposures. Carriers account for 41\%, 38\% and 33\% of
total trade receivables at the three successive dates. At December the note
also identifies two customers individually accounting for at least 10\%, with
shares of 18\% and 11\%. It does not say whether they are carriers
\citep[Note 4, printed p.~38]{apple2024}
\citep[Note 4, printed p.~9]{apple2025q1}. Let $A$ be those two customers,
$B$ the carrier group, and $s(S)$ the share of total trade receivables owed
by customers in a set $S$. Each customer belongs wholly to a set or does not.
Treat the displayed percentages as exact centres for this illustration.
Inclusion-exclusion gives
\[
 s(A\cup B)=s(A)+s(B)-s(A\cap B)
          =29\%+33\%-s(A\cap B).
\]
If both large customers are carriers, their 29\% is already inside the 33\%
carrier share, so the union is 33\%. If neither is a carrier, the union is
62\%. For a concrete construction, allocate receivables across eleven
hypothetical customers with shares
$18,11,9,9,8,7,4,9,9,8,8$ percent. These sum to 100 and only two reach 10\%.
Carriers can be the 18, 11 and 4 customers, totalling 33, or the first
9, 9, 8 and 7 customers, also totalling 33. The same marginal disclosures
therefore permit both endpoint unions. If only one large customer is a carrier,
the union is 44\% or 51\%; exact-centre membership does not permit an
arbitrary continuously divisible overlap. These are conditional percentage
examples, not recovered customer identities or probabilities of loss. The
note's trade denominator is not automatically the balance sheet's \emph{net}
receivable line, so multiplying it by that line would add an unsupported
measurement assumption. Required credit support or collateral applies to
\emph{certain} customers; it does not establish universal recovery.

Vendor non-trade receivables have a separate economic origin and denominator.
Apple buys components from suppliers and sells them to manufacturing vendors
that make subassemblies or final products for Apple. Those component sales
are excluded from product net sales; any gain reduces product cost of sales
when Apple sells the related final products. The two largest reported vendor
shares are 44\% and 23\% at September 2024, and 43\% and 24\% at December,
each pair totalling 67\% on its own dated vendor population
\citep[annual and interim Note 4]{apple2024,apple2025q1}. Equal combined shares
do not establish unchanged counterparties, component cohorts or collection
risk. A funding forecast consequently needs settlements and restrictions on
assets alongside payment dates on obligations. Keeping reported assertions
distinct from those missing facts is the starting point for representing
the company's historical record.
